\documentclass[11pt]{article}

\usepackage[a4paper,margin=1in]{geometry}
\usepackage{amsmath,amssymb,amsfonts,amsthm}
\usepackage{bm}
\usepackage{hyperref}
\usepackage{graphicx}
\usepackage{booktabs}
\usepackage{mathtools}
\usepackage{pgfplots}

\pgfplotsset{compat=1.18}

\hypersetup{
colorlinks=true,
linkcolor=blue,
citecolor=blue,
urlcolor=blue
}

\title{
Selection Principles for the Differentiation Map in the Relational Emergence Model
}

\author{
Yoshifumi Maruko
}

\date{2026}

\begin{document}

\maketitle

\begin{abstract}

The Relational Emergence Model (REM) proposes that relational structure is generated rather than assumed.

A central open problem concerns the selection mechanism of the differentiation map $\mathcal{D}$, which specifies a tensor-product decomposition of Hilbert space.

This paper develops candidate selection principles based on informational and dynamical constraints.

A variational formulation is introduced, together with toy models based on entanglement entropy, mutual information, and energetic stability.

A concrete 3-qubit example shows that informational extremization induces nontrivial relational articulation.

Information geometry, categorical quantum mechanics, and gauge-theoretic subsystem structure are discussed as mathematical frameworks for formalizing differentiation.

These results support the REM thesis that relational structure may emerge through constrained articulation.

\end{abstract}

\section{Introduction}

The Relational Emergence Model (REM) proposes that relational structure is generated rather than assumed.

Differentiation is represented as:

\begin{equation}
\mathcal{D} :
\mathcal{H}
\rightarrow
\mathcal{H}_S \otimes \mathcal{H}_O
\end{equation}

selecting a tensor-product decomposition of Hilbert space.

Hilbert space admits multiple inequivalent factorizations:

\begin{equation}
\mathcal{H}
\cong
\mathcal{H}_A \otimes \mathcal{H}_B
\end{equation}

without a unique preferred structure.

REM distinguishes:

Layer 0:
no preferred factorization

Layer 1:
explicit relational articulation

This paper investigates principles governing this transition.

\section{Variational formulation}

Let:

\begin{equation}
\mathcal{F} = \{F_k\}
\end{equation}

be candidate factorizations.

Define functional:

\begin{equation}
\Phi_k(|\Psi\rangle)
\end{equation}

Selection:

\begin{equation}
k^* =
\arg\operatorname{ext}_k
\Phi_k
\end{equation}

Examples:

\begin{align}
\Phi_S &= S(\rho_A) \\
\Phi_I &= I(A:B) \\
\Phi_H &= -\langle H_{\mathrm{eff}} \rangle
\end{align}

Combined functional:

\begin{equation}
\Phi = \Phi_S + \lambda \Phi_H
\end{equation}

Interpretation of $\lambda$:

\begin{itemize}
\item $\lambda=0$: purely informational articulation
\item small $\lambda$: information-dominated regime
\item large $\lambda$: dynamically stabilized articulation
\end{itemize}

\section{Information geometry of factorization space}

We treat factorization space as a structured manifold:

\begin{equation}
\Phi : \mathcal{F} \rightarrow \mathbb{R}
\end{equation}

Stationary condition:

\begin{equation}
\delta \Phi = 0
\end{equation}

Metric structure:

\begin{equation}
ds^2 =
\sum_{k,l}
g_{kl}(\mathcal{F})
dF^k dF^l
\end{equation}

Possible choices include Fisher metric, Bures metric, or relative-entropy geometry.

Differentiation corresponds to extremal structure on $\mathcal{F}$.

\section{Toy informational model}

Von Neumann entropy:

\begin{equation}
S(\rho)
=
-\mathrm{Tr}(\rho \log_2 \rho)
\end{equation}

Mutual information:

\begin{equation}
I(A:B)
=
S(A)+S(B)-S(AB)
\end{equation}

\section{3-qubit example}

\begin{equation}
|\Psi\rangle
=
\sqrt{0.6}|000\rangle
+
\sqrt{0.3}|111\rangle
+
\sqrt{0.1}|101\rangle
\end{equation}

\begin{center}
\begin{tabular}{lcc}
\toprule
partition & entropy & mutual information \\
\midrule
A vs BC & 0.971 & 1.942 \\
B vs AC & 0.881 & 1.762 \\
C vs AB & 0.971 & 1.942 \\
\bottomrule
\end{tabular}
\end{center}

\section{Category theoretic structure}

Factorizations correspond to monoidal decompositions:

\begin{equation}
X \cong A \otimes B
\end{equation}

Differentiation may be expressed via Frobenius algebra structures:

\begin{equation}
\mu : A \otimes A \rightarrow A
\end{equation}

which define compositional structure selecting relational subsystems.

\section{Gauge theoretic analogy}

Gauge constraints prevent naive tensor decomposition:

\begin{equation}
\mathcal{H}
\not\cong
\mathcal{H}_A \otimes \mathcal{H}_B
\end{equation}

Extended Hilbert space introduces boundary degrees of freedom restoring subsystem factorization.

Differentiation may analogously define relational subsystems consistent with global constraints.

\section{Outlook}

Possible directions include:

\begin{itemize}
\item information-geometric structure of $\mathcal{F}$
\item categorical subsystem emergence
\item tensor-network articulation
\item renormalization of relational layers
\item integration with quantum integrated information
\end{itemize}

Recent work (Zaghi 2025) defines integrated information using quantum Jensen-Shannon divergence minimized over bipartitions.

Such approaches suggest extensions of $\Phi$ beyond entropy and mutual information toward integrated-information criteria.

\section{Conclusion}

Differentiation may be formulated as an optimization problem over Hilbert-space factorizations.

Informational and dynamical constraints provide candidate principles.

REM therefore admits mathematically structured extensions within quantum information theory.

\appendix

\section{Entropy derivation}

\begin{equation}
\rho_A =
\begin{pmatrix}
0.6 & 0 \\
0 & 0.4
\end{pmatrix}
\end{equation}

\begin{equation}
S_A = 0.971
\end{equation}

\begin{equation}
S_B = 0.881
\end{equation}

\begin{equation}
S_C = 0.971
\end{equation}

For pure bipartitions:

\begin{equation}
I(A:\bar A) = 2 S(A)
\end{equation}

confirming consistency between entropy and mutual information criteria.

\section{Mutual information values}

\begin{align}
I(A:BC) &= 1.942 \\
I(B:AC) &= 1.762 \\
I(C:AB) &= 1.942
\end{align}

\begin{thebibliography}{99}
\bibitem{rovelli1996} Rovelli, C. Relational quantum mechanics. Int. J. Theor. Phys. 35, 1637 (1996).
\bibitem{zanardi2001} Zanardi, P. Virtual quantum subsystems. Phys. Rev. Lett. 87, 077901 (2001).
\bibitem{zanardi2004} Zanardi, P., Lidar, D., Lloyd, S. Observable-induced tensor product structures. Phys. Rev. Lett. 92, 060402 (2004).
\bibitem{zurek2003} Zurek, W. Decoherence, einselection, and the quantum origins of classicality. Rev. Mod. Phys. 75, 715 (2003).
\bibitem{carroll2021} Carroll, S., Singh, A. Quantum mereology. Phys. Rev. A 103, 022213 (2021).
\bibitem{abramsky2009} Abramsky, S., Coecke, B. Categorical quantum mechanics. Handbook of Quantum Logic and Quantum Structures (2009).
\bibitem{donnelly2016} Donnelly, W., Freidel, L. Local subsystems in gauge theory and gravity. JHEP 2016.
\bibitem{adlam2023} Adlam, E., Rovelli, C. Information and relational quantum mechanics. Found. Phys. 53 (2023).
\bibitem{riedel2024} Riedel, T. Relational quantum mechanics, quantum relativism, and the iteration of relativity. Stud. Hist. Phil. Sci. 104 (2024).
\bibitem{zaghi2025} Zaghi, A. Integrated Information in Relational Quantum Dynamics. Entropy (2025).
\end{thebibliography}

\end{document}
